\section{Conclusion}

\subsection{What was achieved}

The project delivers a complete, working compiler front end for a language that has no
standard written form. Concretely:

\begin{itemize}[leftmargin=1.4em, itemsep=0.2em]
  \item A hand-written scanner and lexer built on five documented regular expressions,
        resolving \LexiconTotal\ lexicon entries across \OriginCount\ donor languages into
        \TerminalCount\ terminal classes, with multiword recognition, accent folding and
        homograph protection.
  \item A context-free grammar of \ProductionCount\ productions authored in natural
        left-recursive form, mechanically transformed in \LedgerSteps\ recorded steps into
        an \ExecProductionCount-production LL(1) grammar.
  \item FIRST and FOLLOW sets and a parsing table that is verified conflict-free by an
        automated test.
  \item A table-driven predictive parser that returns an explicit derivation on acceptance
        and an expected-terminal set plus concrete repair suggestions on rejection.
  \item A web application that exposes every one of those artefacts for inspection, so the
        compiler is demonstrable rather than merely described.
\end{itemize}

\subsection{What we learned}

The most useful lesson was that the transformations taught in the course are not
bookkeeping. Eliminating left recursion on \nonterm{Seq} \emph{created} a new common prefix
on \term{SEP} that had to be factored out afterwards --- the two techniques interact, and the
order they are applied in decides whether the result is LL(1). Seeing that happen on our own
grammar, in the ledger the tool recorded, taught more than the worked examples did.

The second lesson concerns honesty in tooling. It was tempting, each time a statement was
rejected, to add a production until it passed. We resisted this, and the two remaining
rejections in Section~\ref{sec:results} are documented with the specific LL(1) conflict that
prevents a fix. A grammar that accepts everything distinguishes nothing.

\subsection{Remaining work}

\begin{todobox}[Outstanding before submission]
\begin{enumerate}[leftmargin=1.4em, itemsep=0.2em]
  \item \textbf{Field data collection.} 10--15 real statements must be collected,
        transcribed by hand, cross-checked by a second member, and entered with
        \texttt{source\_kind = field} and the attestation recorded. This is the single
        outstanding graded component.
  \item \textbf{Re-export.} Run \texttt{python tools/export\_tables.py} afterwards; all
        data tables in this report regenerate from the real corpus.
  \item \textbf{Screenshots.} Replace the placeholders in Appendix~\ref{app:screens} with
        captures of the running analyzer.
  \item \textbf{Presentation.} Prepare the ten-minute slide deck and live demonstration.
  \item \textbf{Deployment hardening.} Replace the placeholder signup code before the
        public instance is exposed.
\end{enumerate}
\end{todobox}

\subsection{Future extensions}

Splitting \term{PRON} into full pronouns and clitics would remove both documented
rejections at the cost of re-annotating the pronoun entries. Beyond that, the natural next
step is a semantic layer: the lexicon already carries glosses and donor languages, so a
gloss-composition pass over the parse tree would turn the analyzer from a recogniser into a
translator --- which is what would make it useful outside a coursework setting.
